% Propietario conservado: /Users/ruben/Documents/New project/output/REVISION_ARGUMENTAL_APERTURAS_CIERRES_20260915/III/spanish_source/manuscrito/sections/03b_vacancias_prefactor.tex
% Ampliación propia de III; no modifica el núcleo común copiado de I.
% Residencia propuesta: antes de Retorno de hoja, incidencia areal--volumétrica.
% Procedencia y controles: VACANCIAS_PREFACTOR_PROCEDENCIA.md, en esta carpeta.
\section{Discrete polarization and the tritic mode of arithmetic projections}
\label{iii:iii:sec:vacancias-prefactor}

Refinement and comparison between the APP sheets provide two
complementary observations. The first follows the accumulated discrepancy
between the number of transitions and the capacity acquired; the second preserves
the mode of the tritic selector in the pairing of the additive and
multiplicative fibres. Discrete polarization and the electronic prefactor
are obtained from these respective operations. Their articulation with the vacuum
requires preserving both the temporal information and the spaces on which
the operators act.

\subsection{Refinement discrepancy and endpoint convention}

The cardinalities of the ternary block and its decimal representation fix
\[
 \rho=\log_{1000}729=\log_{10}9,\qquad
 \delta=1-\rho=\log_{10}(10/9).
\]
We have $0<\delta<1$. Moreover, $\delta$ is irrational: an equality
$\delta=p/q$, with $q>0$, would imply $10^{q-p}=9^q$, which is incompatible
with prime factorization. This mismatch already enters the
encoding of refinement. Here its centred readout is made explicit,
preserving the phase origin and the endpoint convention.

\begin{definition}[Lower and upper encodings]
For $\theta\in\mathbb R$ and $n\in\mathbb Z$, define
\begin{align}
 z_n^{\downarrow}(\theta)
 &=\lfloor(n+1)\delta+\theta\rfloor-\lfloor n\delta+\theta\rfloor,
 \label{iii:iii:vp:lower}\\
 z_n^{\uparrow}(\theta)
 &=\lceil(n+1)\delta+\theta\rceil-\lceil n\delta+\theta\rceil.
 \label{iii:iii:vp:upper}
\end{align}
Both take values in $\{0,1\}$. The arrow distinguishes only the
endpoint convention for the intervals; tritic orientation retains
its earlier definition.
\end{definition}

The lower encoding coincides with the one used in the development
of refinement. The upper encoding recovers the formulation of
polarization using the ceiling function. The two expressions are related
by an exactly determined boundary term.

\begin{proposition}[Change of convention and boundary term]
Let $b_n(\theta)=\mathbf1_{\mathbb Z}(n\delta+\theta)$. Then
\begin{equation}
 z_n^{\uparrow}(\theta)-z_n^{\downarrow}(\theta)
 =b_n(\theta)-b_{n+1}(\theta).
 \label{iii:iii:vp:endpoint}
\end{equation}
In particular, for $N\geq0$,
\[
 \sum_{n=0}^{N-1}(z_n^{\uparrow}-z_n^{\downarrow})
 =b_0-b_N.
\]
\end{proposition}
\begin{proof}
For every real $x$, $\lceil x\rceil=\lfloor x\rfloor+1-
\mathbf1_{\mathbb Z}(x)$. Applying this identity to both endpoints
of each increment gives \eqref{iii:iii:vp:endpoint}; its sum telescopes.
The irrationality of $\delta$ implies that, for fixed $\theta$, at most
one integer index satisfies $n\delta+\theta\in\mathbb Z$.
\end{proof}

With $\theta=0$, the first upper symbol is $1$ and the lower one is $0$;
for $n\geq1$ they coincide. The initial term belongs to the register and is
preserved when the convention is changed. The densities and discrepancy
bounds are transported with that term, without modifying the TPK update
or restarting the history at each phase return.

\subsection{Bounded polarization and exact temporal balance}

\begin{definition}[Centred discrete polarization]
The polarization of the upper encoding, relative to the density
$\delta$, is
\begin{equation}
 P_N^{\uparrow}(\theta)
 =\sum_{n=0}^{N-1}\bigl(z_n^{\uparrow}(\theta)-\delta\bigr),
 \qquad P_0^{\uparrow}=0.
 \label{iii:iii:vp:polarization}
\end{equation}
Its temporal increment is
\begin{equation}
 J_N^{\uparrow}(\theta)
 =P_{N+1}^{\uparrow}(\theta)-P_N^{\uparrow}(\theta)
 =z_N^{\uparrow}(\theta)-\delta.
 \label{iii:iii:vp:current}
\end{equation}
\end{definition}

Centring separates uniform growth of density $\delta$ from
phase information. Thus, $P_N^{\uparrow}$ measures an accumulated imbalance,
not the total number of vacancies. The increment $J_N^{\uparrow}$ preserves
the local pulse and the background relative to which it is evaluated.

\begin{theorem}[Uniform bound and discrete temporal continuity]
For $N\geq0$ and every phase $\theta$,
\begin{equation}
 P_N^{\uparrow}(\theta)
 =\lceil N\delta+\theta\rceil-\lceil\theta\rceil-N\delta,
 \qquad |P_N^{\uparrow}(\theta)|<1.
 \label{iii:iii:vp:bound}
\end{equation}
For any integers $0\leq a<b$,
\begin{equation}
 \sum_{n=a}^{b-1}J_n^{\uparrow}(\theta)
 =P_b^{\uparrow}(\theta)-P_a^{\uparrow}(\theta),
 \qquad
 \left|\sum_{n=a}^{b-1}J_n^{\uparrow}(\theta)\right|<1.
 \label{iii:iii:vp:continuity}
\end{equation}
The density of the vacancy sequence is $\delta$.
\end{theorem}
\begin{proof}
Summing the ceiling increments in \eqref{iii:iii:vp:upper} yields the
first identity. Writing $r(x)=\lceil x\rceil-x\in[0,1)$
gives $P_N^{\uparrow}=r(N\delta+\theta)-r(\theta)$; its absolute
value is less than one. On the interval $[a,b)$, one similarly obtains
$r(b\delta+\theta)-r(a\delta+\theta)$, proving
the second bound, which is sharper than the sum of two individual bounds.
Finally,
$N^{-1}\sum_{n=0}^{N-1}z_n^{\uparrow}
=\delta+N^{-1}P_N^{\uparrow}\to\delta$.
\end{proof}

The lower encoding admits the same proofs with
$r(x)$ replaced by $\lfloor x\rfloor-x\in(-1,0]$. In particular,
$P_N^{\uparrow}-P_N^{\downarrow}=b_0-b_N$. The bounded discrepancy
previously used in the Dirichlet readouts and the present polarization
therefore preserve an explicit relationship of origin and endpoints.
Nor does aperiodicity introduce secular drift in this imbalance:
the history can extend indefinitely while the centred discrepancy
remains uniformly bounded.

\subsection{Realization on the charge and time lines}

To realize the balance dimensionally, specify a charge line
$\mathcal L_Q$ with a nonzero section $u_Q$ and an increasing sequence of
times $t_N$ on an oriented time line. With
$\Delta t_N=t_{N+1}-t_N>0$, the centred charge and the mean current of
each interval are defined by
\begin{equation}
 Q_N=P_N^{\uparrow}u_Q,\qquad
 I_N=\frac{Q_{N+1}-Q_N}{\Delta t_N}
     =\frac{u_Q}{\Delta t_N}\bigl(z_N^{\uparrow}-\delta\bigr).
 \label{iii:iii:vp:dimensional}
\end{equation}
Thus, $Q_N\in\mathcal L_Q$ and
$I_N\in\mathcal L_Q\otimes\mathcal L_T^{-1}$. The section $u_Q$ and
the time intervals are data of this realization; the preceding arithmetic
construction has already determined $P_N^{\uparrow}$ and
$J_N^{\uparrow}$.

\begin{proposition}[Preservation of balance under dimensional realization]
For $a<b$,
\begin{equation}
 Q_b-Q_a=\sum_{n=a}^{b-1}\Delta t_n I_n,
 \qquad
 \left|\frac{Q_b-Q_a}{u_Q}\right|<1.
 \label{iii:iii:vp:dimensional-balance}
\end{equation}
If all intervals have length $t_0$, the two possible currents
are $-\delta u_Q/t_0$ and $(1-\delta)u_Q/t_0$.
\end{proposition}
\begin{proof}
Multiplying each identity in \eqref{iii:iii:vp:dimensional} by
$\Delta t_n$ and summing telescopes the charges. The bound is
\eqref{iii:iii:vp:continuity}; the two local values follow from
$z_n^{\uparrow}\in\{0,1\}$.
\end{proof}

The continuity proved is temporal and refers to the centred charge.
A spatial realization also incorporates the localization of charges,
incidence between cells, and identification of boundary fluxes.
In a constitutive description by channels, the coupling is specified
through an excitation map
$\iota:\mathcal L_Q\to\mathcal H_{\rm ch}\otimes\mathcal L_Q$
and an observation map
$\ell:\mathcal H_{\rm ch}\otimes\mathcal L_Q\to\mathcal L_Q$,
where $\mathcal H_{\rm ch}$ is a real realization of the channel space.
If both are
linear and independent of the time index, the response of a fixed operator
$\mathcal V$ is transported by
\[
 Q_N^{\rm resp}=\ell(\mathcal V\otimes I_{\mathcal L_Q})\iota(Q_N),
 \qquad
 I_N^{\rm resp}=\frac{Q_{N+1}^{\rm resp}-Q_N^{\rm resp}}{\Delta t_N}.
\]
Linearity preserves the telescoping balance. This composition makes
explicit the data of a constitutive identification: the scalar sequence
determines the pulse, while $\iota$, $\mathcal V$, and $\ell$
determine its excitation, transport, and observation through channels. Identities
\eqref{iii:iii:vp:polarization}--\eqref{iii:iii:vp:dimensional-balance}
establish the first construction independently of that choice of
realization. The vacuum operator is constructed in
Section~\ref{iii:iii:sec:constitutiva} on its own channels and projectors.

\subsection{Pairing the additive and multiplicative fibres}

The second observation comes from the three-dimensional realization of APP,
$\mathcal A_3=D_9^3$, with uniform measure and
$D_9=\{1,\ldots,9\}$. The observables preserved are
\[
 \Sigma(x)=\rho_9(x_1+x_2+x_3),\qquad
 \Pi(x)=\rho_9(x_1x_2x_3),
\]
where $\rho_9(n)=1+[(n-1)\bmod9]$ for $n>0$. In
$\ell^2(\mathcal A_3)$, let $P_{\Sigma,3}$ and $P_{\Pi,3}$ be
the conditional expectations on their fibres. For example,
\[
 (P_{\Sigma,3}f)(x)=
 \frac1{|\Sigma^{-1}(\Sigma(x))|}
 \sum_{y:\,\Sigma(y)=\Sigma(x)}f(y).
\]
They are orthogonal projections defined on the same support. The incidence
matrix $N_{ab}=|\Sigma^{-1}(a)\cap\Pi^{-1}(b)|$ results from
enumerating the $729$ triples, adding and multiplying their coordinates before
applying $\rho_9$:
\begin{equation}
 N=9\begin{pmatrix}
 0&1&1&0&1&1&0&1&4\\
 1&0&1&1&0&1&1&0&4\\
 1&0&2&0&1&2&0&0&3\\
 0&1&1&0&1&1&0&1&4\\
 1&0&1&1&0&1&1&0&4\\
 0&0&2&1&0&2&0&1&3\\
 0&1&1&0&1&1&0&1&4\\
 1&0&1&1&0&1&1&0&4\\
 0&1&2&0&0&2&1&0&3
 \end{pmatrix}.
 \label{iii:iii:vp:joint-matrix}
\end{equation}
The rows sum to $81$, and the column sums are
\[
 (m_1,\ldots,m_9)=(36,36,108,36,36,108,36,36,297).
\]
Therefore, the pairing of the normalized indicator functions of the fibres
is $C_{ab}=N_{ab}/\sqrt{81m_b}$. Normalization incorporates the
multiplicities produced by APP, not a choice of singular values.

\begin{proposition}[Singular mode of the tritic selector]
In the canonical chart, the selector vector is
$m_{\rm ph}=(1,-1,0)^3$. We have
\begin{equation}
 Cm_{\rm ph}=C^{\mathsf T}m_{\rm ph}=-\tfrac12m_{\rm ph}.
 \label{iii:iii:vp:trit-mode}
\end{equation}
Thus, the corresponding singular value is $s=1/2$.
\end{proposition}
\begin{proof}
In the six columns where $m_{\rm ph}$ is nonzero, $m_b=36$ and
$C_{ab}=N_{ab}/54$. Multiplying the rows of
\eqref{iii:iii:vp:joint-matrix} by $(1,-1,0)^3$ yields
$-m_{\rm ph}/2$. For the transposed product, the signed sums
of columns $3,6,9$ are zero; the other six give the same
values $-m_{{\rm ph},b}/2$. This proves both identities and, by
composition, $CC^{\mathsf T}m_{\rm ph}=m_{\rm ph}/4$.
\end{proof}

\subsection{Commutator norm and electronic prefactor}

\begin{theorem}[Exact norm of arithmetic incompatibility]
The two projections satisfy
\begin{equation}
 \bigl\|[P_{\Sigma,3},P_{\Pi,3}]\bigr\|=\frac{\sqrt3}{4}.
 \label{iii:iii:vp:prefactor}
\end{equation}
The maximum is attained in the principal plane corresponding to the
tritic mode of \eqref{iii:iii:vp:trit-mode}.
\end{theorem}
\begin{proof}
The rational matrix $M=CC^{\mathsf T}$ has entries
$M_{ac}=\sum_bN_{ab}N_{cb}/(81m_b)$. Its complete spectrum is
verified through the following independent subspaces: the
vectors $(1,\ldots,1)$, $(1,-1,0)^3$, and $(1,1,-2)^3$ have
eigenvalues $1$, $1/4$, and $7/132$, respectively; the subspace of
vectors supported on $\{3,6,9\}$ whose sum is zero has dimension
$2$ and eigenvalue $1/18$; the zero-sum vectors supported on
$\{1,4,7\}$ or $\{2,5,8\}$ form a kernel of dimension $4$.
Substitution into the matrix verifies these actions, and the dimensions
sum to $9$. Consequently,
\[
 \operatorname{spec}M=
 \{1,1/4,7/132,1/18,1/18,0,0,0,0\}.
\]
For two orthogonal projections, a singular value $s\in(0,1)$ of the
pairing determines a principal plane in which their matrices are
\[
 P=\begin{pmatrix}1&0\\0&0\end{pmatrix},\qquad
 Q=\begin{pmatrix}s^2&s\sqrt{1-s^2}\\
 s\sqrt{1-s^2}&1-s^2\end{pmatrix}.
\]
This form is obtained by taking a unit vector $u$ in the first
image with $PQu=s^2u$ and completing it with
$(Qu-s^2u)/(s\sqrt{1-s^2})$. Their commutator is
$s\sqrt{1-s^2}\bigl(\begin{smallmatrix}0&1\\-1&0\end{smallmatrix}\bigr)$.
The sectors with singular value $0$ or $1$ contribute zero.
Therefore, the total norm is the maximum of
$\sqrt{\lambda(1-\lambda)}$ over the spectrum of $M$. That maximum
is $\sqrt{3/16}$, attained at $\lambda=1/4$.
\end{proof}

The norm is the prefactor entering the electronic composition.
Its derivation preserves the three-dimensional support, the fibres, and the
selector mode; scalar evaluation does not replace that structure. In
particular, knowing $\sqrt3/4$ gives the maximum magnitude of
the incompatibility, while the projectors preserve the directions
in which it acts.

\clearpage
\subsection{Quadratic defect and response of the \(90/120\) sectors}

The tritic mode simultaneously determines
\begin{equation}
 1-s^2=\frac34=\frac{90}{120},\qquad s=\frac12.
 \label{iii:iii:vp:defect-ratio}
\end{equation}
The first expression is the quadratic defect of the APP principal plane.
The last preserves the lengths of the two calendar sectors.
The constitutive response uses those sectors on another object: a
contractive channel $q\in(0,1)$. Its evaluation is
\[
 r(q)=\frac{1-q^{90}}{1-q^{120}}.
\]
The relationship between the sector ratio and this response can be
made precise without identifying their operators.

\begin{proposition}[Limit of the response ratio]
Para $0<q<1$,
\begin{equation}
 \frac34<r(q)<1,\qquad
 \lim_{q\uparrow1}r(q)=\frac34=1-s^2.
 \label{iii:iii:vp:response-limit}
\end{equation}
\end{proposition}
\begin{proof}
With $u=q^{30}\in(0,1)$, factoring $1-u^3$ and
$1-u^4$ gives
$r(q)=(1+u+u^2)/(1+u+u^2+u^3)$. The limit at $u=1$ is
$3/4$. Moreover,
$4(1+u+u^2)-3(1+u+u^2+u^3)
=(1-u)(3u^2+2u+1)>0$, and the denominator exceeds the numerator
by $u^3>0$. This yields both inequalities.
\end{proof}

The exact coincidence concerns the defect of the singular mode and the limit
of the sector response. In the contractive interior, the response
depends on $q$ and is strictly greater than $3/4$. The constitutive
development preserves its two channels $q_\pm$ and its projectors
$P_\pm$; the projectors $P_{\Sigma,3}$ and $P_{\Pi,3}$, for their
part, preserve the arithmetic fibres that determined the prefactor.
An operator equivalence between the two realizations requires a map
between their spaces that intertwines the corresponding actions.
Equations \eqref{iii:iii:vp:trit-mode}, \eqref{iii:iii:vp:prefactor}, and
\eqref{iii:iii:vp:response-limit} specify the shared information and
keep the data of each realization visible.
